\begin{table}[ht]
\caption{\textbf{Example of Gemini reasoning on \sustainable, continue}}
\centering
% \resizebox{\textwidth}{!}{

\begin{tabular}{p{13cm}}
\toprule
\textit{\textbf{Example positive document}} \\
\midrule
Soluble Salts

Soluble salts may accumulate on the top of the soil, forming a yellow or white crust. A ring of salt deposits may form around the pot at the soil line or around the drainage hole. Salts may also build up on the outside of clay pots. In house plants, signs of excess soluble salts include reduced growth, brown leaf tips, dropping of lower leaves, small new growth, dead root tips, and wilting.

Soluble salts are minerals dissolved in water. Fertilizer dissolved in water becomes a soluble salt. When water evaporates from the soil, the minerals or salts stay behind. As the salts in the soil become more and more concentrated, it becomes more difficult for plants to take up water. If salts build up to an extremely high level, water can be taken out of the root tips, causing them to die. High levels of soluble salts damage the roots directly, weakening the plant and making it more susceptible to attack from insects and diseases. One of the most common problems associated with high salt levels is root rot.

The best way to prevent soluble salt injury is to stop the salts from building up. When watering, allow some water to drain through the container and then empty the saucer. Do not allow the pot to sit in water. If the drained water is absorbed by the soil, the salts that were washed out are reabsorbed through the drainage hole or directly through a clay pot. 

...\\

\bottomrule
\end{tabular}
% }
\label{tab:example_sustainable_living_reason1}
\end{table}